% ************************************************************
% CONCLUSION
% ************************************************************

\section{Conclusion}

This project develops a reproducible framework for testing persistent homology as an auxiliary diagnostic for symmetric-cipher output randomness. In the current 64-replicate internal run, byte-pair H0 summaries strongly separated the LCG weak-control condition from the SHA-256 seeded baseline, while AES-CTR deterministic-output conditions, OS CSPRNG output, and xorshift32 did not show comparable separation under the tested embeddings and feature set. The study is deliberately scoped as computational diagnostics rather than cryptanalytic attack development. Its publication value rests on transparent data generation, careful controls, reproducible code, evidence-register traceability, and restrained interpretation. The next evidentiary step is to import external randomness-test results and evaluate whether those conventional diagnostics agree with, complement, or fail to detect the same weak-control structure observed in the topological summaries.
